In the strategy comparison, Proposed-approach achieved macro-$F_1$ scores of 88.25\% and 80.68\% on Internal and Cao2018, respectively, exceeding both baselines on every dataset (Table~\ref{tab:f1_significance}). Two-sided Wilcoxon signed-rank tests on session-level $F_1$, Bonferroni-corrected across the four baseline comparisons in that table (two baselines $\times$ two datasets), confirm that the advantage over both BLINKER and MNE is significant on both corpora ($p<0.05$ after correction in every case). This performance gap is reflected in the precision-recall operating points shown in Figure~\ref{fig:exp2_pr_scatter}, where BLINKER, the stronger of the two baselines, consistently favors recall over precision on both corpora. This asymmetry follows from the pooled-threshold mechanism described earlier. Deriving a single channel threshold from a mixture of blink-heavy and blink-free epochs lowers the effective detection threshold relative to an epoch-aware estimate, so more candidate events cross it, raising recall but also admitting more false positives that lower precision.